\section{Affine Dual: Coordinate Matroid}\label{sec:affine}
%==============================================================================


The preceding sections developed the main graph-capacity and equality arc for partial views. Computing the resulting graph quantities is hard in general. This section isolates a regime in which the coordinate side of the model becomes tractable: when the realized state family is affine, semantic determination reduces to standard linear algebra, and the resulting upper bounds can be computed by matrix rank rather than by solving a general graph-capacity problem.

The resulting object is a representable matroid on fact indices. It is the coordinate-side dual of the same realized-state structure that generates the confusability graph, and it provides computationally tractable upper bounds on confusability and capacity. The tractability claim is representation-sensitive: throughout this section, the affine family $A=a_0+V$ is assumed to be given by an explicit linear presentation of $V$, for example a basis or generator matrix over $\mathbb F_q$, together with the explicit coordinate-subset view family. Under that input model, the relevant quantity is just the rank of a restricted coordinate map and is computable by Gaussian elimination.

\paragraph{Warm-up example.}
Before stating the general proposition, consider the binary square from Section~\ref{sec:binary-square}. Here the realized state family is the full affine space $A=\mathbb F_2^2$, so we may take $a_0=(0,0)$ and direction space $V=\mathbb F_2^2$ with generator matrix
\[
G = \begin{bmatrix}1 & 0 \\\ 0 & 1\end{bmatrix}.
\]
The coordinate functionals are $e_1=[1\ \ 0]$ and $e_2=[0\ \ 1]$. A single-coordinate view, say $S=\{1\}$, has rank $t(S)=1$: it captures one of the two available dimensions, so it reveals exactly half of the informational dimension of the state space. This is the finite-affine meaning of the later matroid rank bound: the coordinate-side quantity $t(S)$ measures how much of the state survives the view, and in the affine regime it is obtained by ordinary matrix rank.

\begin{proposition}[Affine fact-matroid specialization]\label{thm:affine-fact-matroid}
Assume the realized state family is affine over a field, so that the valid latent tuples form an affine translate of a linear subspace of the ambient fact space. For any fact set $S$ and fact index $i$, semantic determination of $i$ by $S$ is equivalent to membership of the coordinate functional for $i$ in the linear span of the coordinate functionals indexed by $S$. Consequently the fact indices carry a representable matroid: a fact set is independent exactly when its coordinate functionals are linearly independent, and the minimal determining fact sets are exactly the bases, all of common cardinality equal to the rank.\leanmeta{\LHrng{AFM}{1}{5}}
\end{proposition}

\begin{proof}
Let the affine realized-state family be $A=a_0+V$, where $a_0$ is a fixed origin and $V$ is a linear subspace of the ambient fact space. Fix a fact set $S$ and a fact index $i$.

By definition, fact $i$ is semantically determined by $S$ exactly when for all $x,x'\in A$, agreement on every coordinate in $S$ forces agreement on coordinate $i$. Writing $d=x-x'$, this is equivalent to the statement that every direction vector $d\in V$ whose coordinates in $S$ vanish also satisfies $d_i=0$.

Let $e_j$ denote the coordinate functional for fact $j$. The condition above says precisely that
\[
V\cap \bigcap_{j\in S}\ker(e_j)\subseteq \ker(e_i).
\]
In finite-dimensional linear algebra, this is equivalent to $e_i$ lying in the span of the coordinate functionals $\{e_j:j\in S\}$. Thus semantic determination by $S$ is exactly span membership of the corresponding coordinate functional.

The coordinate functionals therefore realize a representable matroid on fact indices. In that matroid, independence is linear independence of the coordinate functionals, bases are the maximal independent spanning sets, and every basis has the common rank cardinality. Since span corresponds exactly to semantic determination, the minimal determining fact sets are precisely those bases.
\end{proof}

The matroid and the confusability graph are complementary objects induced by the same realized-state model. Under the common specialization to finite affine families with coordinate-projection views, the matroid rank provides \emph{computationally tractable} upper bounds on confusability and capacity. For each coordinate set $S$ the quantity $t(S)$ is the rank of the restricted coordinate map $\pi_S|_V$.\leanmeta{\LHrng{AFM}{10}{11}} The formalization proves the basic size bounds $t(S)\le |S|$ and $t(S)\le \dim V$, and shows that $t(S)=|S|$ on linearly independent coordinate families while $t(S)$ reaches the full ambient determining rank exactly on determining sets.\leanmeta{\LHrng{AFM}{6}{9}, \LH{AFM12}} Hence each $t(S)$ is computable by Gaussian elimination on an explicit presentation of $V$, and the bounds of Corollary~\ref{cor:matroid-capacity-bounds} are computable in polynomial time. The utility is immediate: $t(S)$ provides an explicit upper bound on the Shannon capacity of the induced confusability graph, bypassing the hardness of the general capacity computation. Computing Shannon capacity for general graphs is hard (the independence number is NP-hard, and even the Lovász-$\vartheta$ bound requires semidefinite programming~\cite{pekosz2011complexity,grotschel1981ellipsoid}). In contrast, once the affine family is explicitly presented, the relevant rank quantities reduce to standard linear algebra.

\subsection{View-fiber dimensions and capacity}
\label{sec:affine-to-graph}

Let $A=a_0+V$ with $V$ an $r$-dimensional linear subspace over a finite field $\mathbb{F}_q$, and let admissible views be coordinate-subset projections. For a coordinate set $S$ write $t(S)$ for the rank of the corresponding coordinate functionals on $V$.

\leanmetapending{\LH{MFT86}}
\begin{proposition}[View-fiber clique sizes]
Let $A$ and $t(S)$ be as above. Each projection-fiber for view $S$ is an affine subspace of dimension $r-t(S)$ and size $q^{r-t(S)}$. The single-view confusability graph $G_S$ is therefore a disjoint union of $q^{t(S)}$ cliques each of size $q^{r-t(S)}$. In particular
\[
\alpha(G_S)=q^{t(S)},\qquad \Theta(G_S)=t(S)\log q,
\]
where $\alpha$ is the independence number and $\Theta$ is Shannon capacity (logarithms in the same base used elsewhere).
\end{proposition}

\begin{proof}
The projection onto coordinates $S$ is an affine-linear map whose linear part restricted to $V$ has rank $t(S)$. Its kernel is therefore a linear subspace of $V$ of dimension $r-t(S)$, so every fiber is an affine translate of that kernel and has cardinality $q^{r-t(S)}$. The fibers partition $A$, and within each fiber every pair of states is confusable under $S$, so $G_S$ is a disjoint union of equal-size cliques. The number of fibers equals $|A|/q^{r-t(S)}=q^{t(S)}$, so one can select one representative per fiber to obtain an independent set of size $q^{t(S)}$, whence $\alpha(G_S)=q^{t(S)}$. The block-power identity for such clustered graphs gives $\alpha(G_S^{\boxtimes n})=q^{nt(S)}$, and normalization yields $\Theta(G_S)=t(S)\log q$.
\end{proof}

\leanmetapending{\LH{FM6}, \LHrng{AFM}{6}{12}}
\begin{corollary}[Matroid ranks give capacity bounds]\label{cor:matroid-capacity-bounds}
Let $G$ be the full confusability graph generated by a family of coordinate-subset views $\mathcal S$. Then
\[
\alpha(G)\le \min_{S\in\mathcal S} q^{t(S)}
\qquad\text{and}\qquad
\Theta(G) \le \min_{S\in\mathcal S} t(S)\log q.
\]
Thus the matroid rank function $t(\cdot)$ on coordinate sets supplies explicit upper bounds on one-shot and asymptotic exact-recovery rates for the induced confusability graph.
\end{corollary}

\begin{proof}
An independent set for $G$ must be independent in each single-view graph $G_S$, hence its size is at most $\alpha(G_S)=q^{t(S)}$ for every $S\in\mathcal S$. The inequalities follow from taking the minimum over $\mathcal S$ and applying the block-power argument from the proposition.
\end{proof}

\noindent\textbf{Remark.} The matroid rank $t(S)$ is the ``informational dimension'' captured by coordinates in $S$, directly analogous to the rank function in representable matroids studied in combinatorial optimization~\cite{oxley2011matroid,welsh1976matroid,recski1989matroid} and coding theory~\cite{cover2006elements}. Here $t(S)=r$ iff $S$ determines the entire state, and any admissible view containing a basis removes confusability entirely (the induced graph is edgeless). These finite-affine, coordinate-view statements show the matroid is directly applicable to the graph/capacity arc; extensions beyond this specialization are directions for future work.

\paragraph{Running example: Binary square.}
The binary square $\{0,1\}^2$ from Section~\ref{sec:binary-square} is an affine space over $\mathbb{F}_2$ with $r=2$. Each single-coordinate view has matroid rank $t(\{1\}) = t(\{2\}) = 1$, giving fiber size $2^{2-1} = 2$ and single-view capacity bound $\Theta(G_{\{i\}}) \le 1 \cdot \log 2 = \log 2$. The full confusability graph (the 4-cycle) achieves this bound, so the matroid rank bound is tight in this case. The matroid viewpoint shows that the 4-cycle structure arises precisely because each view captures only half the informational dimension.
